\chapter*{Preface}
\addcontentsline{toc}{chapter}{Preface}
\markboth{Preface}{Preface}\label{ch:preface}

This book documents a complete bare-metal stack for the dual-core ESP32-S3,
written almost entirely in Ada and running with no operating system, no FreeRTOS,
and no ESP-IDF underneath it. The application boots through a 2nd-stage bootloader
written for the purpose, comes up on both cores under a from-scratch Ada runtime,
and drives every peripheral through hand-written drivers. From the GPIO blink to a
pure-Ada TLS~1.3 client fetching live data over HTTPS, the code in these pages runs
on real silicon.

\section*{The thesis}

The whole project asks one question: how far can bare-metal Ada go on its own? The
usual answer on a microcontroller is ``not far without a C library,'' so the usual
embedded stack reaches for the vendor SDK, links a C TLS library, and trusts a
pre-built runtime. This project does the opposite. It builds the bootloader, the
tasking runtime, the peripheral drivers, the filesystem, the network stack, and the
cryptography in Ada, and it leans on the language to make that safe: range-checked
types catch a wrong-width register write at compile time, a private type guards
every hardware resource, and the Jorvik tasking profile keeps the whole program
statically analysable. The payoff is a system where every byte of flash and RAM,
and the worst-case depth of every stack, can be accounted for before the board ever
runs.

\section*{Who this book is for}

The reader we wrote for is an experienced \textbf{C or systems programmer} who is
new to Ada, to bare-metal work, or to both. You are assumed to know what a stack, a
cache, an interrupt and a peripheral register are. You are \emph{not} assumed to
know Ada: Chapter~\ref{ch:ada-primer} is a primer for exactly this reader, and the
recurring ``\textbf{From C / ESP-IDF}'' asides translate each idea back to ground you
already stand on. An Ada programmer curious about bare metal, or an embedded
engineer evaluating Ada against C, will also find the parts they need.

\section*{What you need}

One \textbf{ESP32-S3 board} with the built-in USB-Serial-JTAG port (most devkits)
and one USB cable is enough to run the core of the book. A few later chapters add
external parts (an SD card, a W5500 Ethernet module, a display, a sensor); each
example names its own wiring. The toolchain comes from \textbf{Alire}, the Ada
package manager, which supplies the cross-compiler and pins the runtime crate. No
other SDK is installed. Chapter~\ref{ch:getting-started} takes you from a fresh
clone to a blinking LED in a few commands.

\section*{How the book is organised}

The five parts move from the language outward to the silicon:

\begin{description}[style=nextline, leftmargin=0pt]
\item[Part I, Ada on Bare Metal] the language and the project: a C-programmer's
  Ada primer, the package system that holds the design together, the build, and the
  anatomy of an image.
\item[Part II, The Runtime] what an Ada runtime is, how the tasking half (GNARL)
  and the sequential half fit on bare metal, and what the restricted profiles give
  up to stay analysable.
\item[Part III, Talking to Hardware] the encapsulation patterns the drivers share,
  the peripheral HAL, the tasking and synchronisation model, and putting the chip
  to sleep.
\item[Part IV, Building Real Systems] the storage, networking, Modbus, and TLS
  built on top (the language interpreters --- LISP, Lua, MicroPython, Forth ---
  live in the companion examples collection).
\item[Part V, Under the Hood] the runtime build, the bootloader and PSRAM bring-up,
  production concerns, the heap and the static-memory accounting, context switching,
  interrupts, how the project is tested, and the conformance sweep on real silicon.
\end{description}

\section*{How to read it}

Read Part~I first; after that the chapters are fairly independent, so you can follow
the contents or jump to the subsystem you need. Every concept is anchored in code
that ships in the repository, and almost every chapter points at a runnable example.
The \texttt{./x} command drives all of them: \texttt{./x run <example>} builds,
flashes, and opens the serial console. The code in these pages is not pseudocode.
Every worked listing is compiled against the real HAL as part of the build
(Chapter~\ref{ch:testing}), and the examples are run on hardware, so what you read
is what runs.
